\documentclass[../../main.tex]{subfiles}% 注意这里的文件路径不能用 ./main.tex ,否则用latexmk编译子文件会报错
\graphicspath{{\subfix{./image/}}{\subfix{../../image/}}} % 指定图片引用路径,后续可以直接使用图片文件名
% 如果图片存放在上级目录,则使用 ../../image/ 作为备用路径

% 例如:
% \begin{figure}[H]
% \centering
% \includegraphics[width=0.3\textwidth]{图.png}
% \caption{}
% \label{figure:图}
% \end{figure}
% 注意:上述\label{}一定要放在\caption{}之后,否则引用图片序号会只会显示??.

\begin{document}

\section{一般域上的Mittag-Leffler定理、Weierstrass因子分解定理和插值定理}

\begin{definition}
设 \(D\) 是域, \(f\) 在 \(D\) 上除了极点外, 在其他点处都全纯, 则称 \(f\) 是 \(D\) 上的\textbf{亚纯函数}.
\end{definition}

\begin{theorem}[Mittag-Leffler定理]\label{theorem:Mittag-Leffler定理}
设 \(D\) 是 \(\mathbb{C}\) 中的域, \(\{a_n\}\) 是 \(D\) 中的一列互不相同并且在 \(D\) 内部无极限点的点列, 且 \(\psi_n(z)\) 是给定的一列有理函数, 则必存在 \(D\) 上的亚纯函数 \(f\), 使得 \(f\) 恰以 \(\{a_n\}\) 为其极点集, 并且在每个 \(a_n\) 处的 Laurent 展开式的主要部分恰为 \(\psi_n(z)\).
\end{theorem}
\begin{proof}
对每个 \(a_n\), 存在圆盘 \(B(a_n,3\varepsilon_n)\subset D\), 使得 \(\{B(a_n,3\varepsilon_n)\}\) 互不相交. 由引理 3.7.2, 存在 \(\varphi_n \in C^\infty(\mathbb{C})\), 使得 \(\varphi_n(z)\) 在 \(B(a_n,\varepsilon_n)\) 上恒等于 \(1\), 并且 \(\operatorname{supp}\varphi_n \subset B(a_n,2\varepsilon_n)\). 令
\begin{align*}
u(z) = \sum\limits_{n=1}^\infty \varphi_n(z)\psi_n(z).
\end{align*}
则 \(u \in C^\infty(D \setminus \{a_n: n \in \mathbb{N}\})\), 并且当 \(z \in B(a_n,\varepsilon_n) \setminus \{a_n\}\) 时, 有 \(u(z) = \varphi_n(z)\psi_n(z) = \psi_n(z)\). 再令
\begin{align}\label{eq::::AQAIijij-2}
h(z) = \begin{cases}
\frac{\partial u(z)}{\partial \overline{z}}, & z \in D \setminus \{a_n: n \in \mathbb{N}\}; \\
0, & z \in \{a_n: n \in \mathbb{N}\},
\end{cases}
\end{align}
注意到当 \(z \in B(a_n,\varepsilon_n) \setminus \{a_n\}\) 时, \(h(z) = \frac{\partial u(z)}{\partial \overline{z}} = \frac{\partial \psi_n(z)}{\partial \overline{z}} = 0\), 由 (3) 式便知 \(h \in C^\infty(D)\). 根据\cref{theorem:复变-定理3.7.3}, 下面的 \(\overline{\partial}\) 问题
\begin{align*}
\frac{\partial v}{\partial \overline{z}} = h
\end{align*}
有解 \(v \in C^\infty(D)\). 令 \(f = u - v\), 则 \(f\) 在 \(D \setminus \{a_n: n \in \mathbb{N}\}\) 上全纯, 这是因为 \(\frac{\partial f}{\partial \overline{z}} = \frac{\partial u}{\partial \overline{z}} - \frac{\partial v}{\partial \overline{z}} = h - h = 0\) 的缘故. 由\eqref{eq::::AQAIijij-2} 式及 \(v \in C^\infty(D)\), 即知 \(f\) 在每个 \(a_n\) 处的 Laurent 展开式的主要部分为 \(\psi_n(z)\).
\end{proof}

\begin{theorem}[Weierstrass因子分解定理]\label{theorem:Weierstrass因子分解定理}
设 \(D\) 是 \(\mathbb{C}\) 中的单连通域, \(\{a_n\}\) 是 \(D\) 中的一列互不相同并且在 \(D\) 内部无极限点的点列, \(\{k_n\}\) 是一列自然数, 则必存在 \(D\) 上的全纯函数 \(f\), 使得 \(f\) 恰以 \(\{a_n\}\) 为其零点集, 并且 \(f\) 在每个 \(a_n\) 处的零点阶数恰为 \(k_n\).
\end{theorem}
\begin{proof}
由 \hyperref[theorem:Mittag-Leffler定理]{Mittag-Leffler定理}, 存在 \(D\) 上的亚纯函数 \(g\), 使得 \(g\) 恰以 \(\{a_n\}\) 为其极点集, 并且 \(g\) 在每个 \(a_n\) 处的 Laurent 展开式的主要部分恰为 \(\frac{k_n}{z-a_n}\). 于是, 在 \(a_n\) 附近有
\begin{align}\label{eq::444jJJ}
g(z) = \frac{k_n}{z-a_n} + g_n(z),
\end{align}
这里 \(g_n(z)\) 在 \(a_n\) 附近全纯. 固定 \(a_0 \in D \setminus \{a_n: n \in \mathbb{N}\}\), 令
\begin{align*}
F(z) = \int_{a_0}^z g(\zeta) \mathrm{d}\zeta, \quad z \in D \setminus \{a_n: n \in \mathbb{N}\}.
\end{align*}
则对于固定的 \(z\), 由于 \(D \setminus \{a_n: n \in \mathbb{N}\}\) 中连接 \(a_0\) 和 \(z\) 的曲线不同, 所得的积分值也不同. 但从\eqref{eq::444jJJ}式可看出, 两个不同的积分值之差是 \(2\pi \mathrm{i}\) 的整数倍. 因此, \(F(z)\) 是 \(D \setminus \{a_n: n \in \mathbb{N}\}\) 上的多值全纯函数, 但 \(F(z)\) 在同一点处的任意两个函数值之差是 \(2\pi \mathrm{i}\) 的整数倍. 于是, \(f(z) = e^{F(z)}\) 便是 \(D \setminus \{a_n: n \in \mathbb{N}\}\) 上的单值全纯函数. 我们可验证 \(f\) 就是要求的 \(D\) 上的全纯函数. 事实上, 由 (4) 式和 (5) 式可知, 在 \(a_n\) 附近有
\begin{align*}
F(z) = k_n \ln(z-a_n) + h_n(z),
\end{align*}
这里 \(h_n(z)\) 在 \(a_n\) 附近全纯, 于是
\begin{align*}
f(z) = (z-a_n)^{k_n} e^{h_n(z)}.
\end{align*}
所以 \(f\) 是 \(D\) 上的全纯函数, 恰以 \(\{a_n: n \in \mathbb{N}\}\) 为其零点集, 并且在每个 \(a_n\) 处的零点阶数为 \(k_n\).
\end{proof}

\begin{theorem}
(插值定理) 设 \(\{a_n\}\) 是单连通域 \(D\) 中的一列互不相同并且在 \(D\) 内部无极限点的点列, 且 \(P_n(z)\) 是给定的一列多项式, 则必存在 \(D\) 上的全纯函数 \(f\), 使得 \(f\) 在每个 \(a_n\) 处的 Taylor 级数的前 \(k_n+1\) 项部分和恰为 \(P_n(z)\).
\end{theorem}
\begin{proof}
由\hyperref[theorem:Weierstrass因子分解定理]{Weierstrass因子分解定理}, 可取 \(g \in H(D)\), 使得 \(g\) 恰以 \(\{a_n\}\) 为其零点集, \(g\) 在每个 \(a_n\) 处的零点阶数恰为 \(k_n+1\). 对每个 \(a_n\), 存在圆盘 \(B(a_n,3\varepsilon_n)\subset D\), 使得 \(\{B(a_n,3\varepsilon_n)\}\) 互不相交. 由\cref{lemma:引理3.7.2}, 存在 \(\varphi_n \in C^\infty(\mathbb{C})\), 使得 \(\varphi_n(z)\) 在 \(B(a_n,\varepsilon_n)\) 上恒等于 \(1\), 并且 \(\operatorname{supp}\varphi_n \subset B(a_n,2\varepsilon_n)\). 令
\begin{align*}
u(z) = \begin{cases}
\varphi_n(z)\frac{P_n(z)}{g(z)}, & z \in B(a_n,3\varepsilon_n) \setminus \{a_n\}, n \in \mathbb{N}; \\
0, & z \in D \setminus \bigcup\limits_{n=1}^\infty (B(a_n,3\varepsilon_n) \setminus \{a_n\}).
\end{cases}
\end{align*}
则 \(u \in C^\infty(D \setminus \{a_n: n \in \mathbb{N}\})\), 并且当 \(z \in B(a_n,\varepsilon_n) \setminus \{a_n\}\) 时, 有 \(u(z) = \varphi_n(z)\frac{P_n(z)}{g(z)} = \frac{P_n(z)}{g(z)}\). 再令
\begin{align*}
h(z) = \begin{cases}
\frac{\partial u(z)}{\partial \overline{z}}, & z \in D \setminus \{a_n: n \in \mathbb{N}\}; \\
0, & z \in \{a_n: n \in \mathbb{N}\}.
\end{cases}
\end{align*}
注意到当 \(z \in B(a_n,\varepsilon_n) \setminus \{a_n\}\) 时, \(h(z) = \frac{\partial u(z)}{\partial \overline{z}} = \frac{\partial}{\partial \overline{z}}\left(\frac{P_n(z)}{g(z)}\right) = 0\), 由 (8) 式便知 \(h \in C^\infty(D)\). 根据\cref{theorem:复变-定理3.7.3}, 下面的 \(\overline{\partial}\) 问题
\begin{align}
\frac{\partial v}{\partial \overline{z}} = h\label{eq::::AQAIijij-7}
\end{align}
有解 \(v \in C^\infty(D)\). 令
\begin{align*}
f(z) = g(z)u(z) - g(z)v(z),
\end{align*}
则 \(f\) 在 \(D \setminus \{a_n: n \in \mathbb{N}\}\) 上全纯, 这是因为 \(\frac{\partial f}{\partial \overline{z}} = g\left( \frac{\partial u}{\partial \overline{z}} - \frac{\partial v}{\partial \overline{z}} \right) = g(h-h) = 0\) 的缘故. 由\eqref{eq::::AQAIijij-7} 式和 \(v \in C^\infty(D)\) 可知 \(f \in H(D)\), 并且在每个 \(a_n\) 附近有
\begin{align*}
f(z) = P_n(z) - g(z)v(z).
\end{align*}
若注意到 \(a_n\) 是 \(g(z)v(z)\) 的至少 \(k_n+1\) 阶零点, 则知 \(f\) 在 \(a_n\) 处的 Taylor 级数的前 \(k_n+1\) 项部分和便是 \(P_n(z)\).
\end{proof}




\end{document}